\PartChapter{Detailed Theory: Prime Packets, Valuation Refinement, and Exact Obstructions}{Prime-Packet Differentials and Radical Divisibility}

\section{The prime-coordinate module}
Let
\[
 a+b=c,
 \qquad
 \gcd(a,b)=1.
\]
The three integers are pairwise coprime. Write
\[
 S=S_a\sqcup S_b\sqcup S_c
 =\{p:p\mid abc\}.
\]
The first object of the calculus is the free prime-packet module
\[
 \Omega_S=\bigoplus_{p\in S}\Z\xi_p.
\]
For a positive integer $n$, define the universal multiplicative differential
\[
 \boxed{
 \mathrm dn
 =
 n\sum_{p\mid n}\frac{v_p(n)}p\,\xi_p.
 }
\]
Every coefficient is an integer because $p\mid n$. The map obeys the exact Leibniz law
\[
 \mathrm d(mn)=m\,\mathrm dn+n\,\mathrm dm.
\]
The verification is local. For each prime $p$,
\[
 v_p(mn)=v_p(m)+v_p(n),
\]
so the $p$-coordinate of the right side is
\[
 mn\frac{v_p(m)+v_p(n)}p.
\]

A scalar packet is an integer vector
\[
 x=(x_p)_{p\in S}\in\Z^S.
\]
Pairing with $x$ gives the arithmetic derivative
\[
 \boxed{
 \mathrm d_x(n)
 =
 n\sum_{p\mid n}\frac{v_p(n)}p x_p.
 }
\]
The Leibniz rule survives scalarization.

Arithmetic derivatives and their connection with $abc$ have an established literature. The present construction is organized here in prime-packet coordinates and later refined into filtered shell objects~\cite{PastenDerivatives,VojtaABC}.

\section{The additive row}
The condition
\[
 \mathrm d_x(a)+\mathrm d_x(b)=\mathrm d_x(c)
\]
is one integer equation. Define
\[
 \lambda_p=
 \begin{cases}
 \dfrac{a v_p(a)}p,&p\mid a,\\[6pt]
 \dfrac{b v_p(b)}p,&p\mid b,\\[6pt]
 -\dfrac{c v_p(c)}p,&p\mid c.
 \end{cases}
\]
Then
\[
 \boxed{\lambda\cdot x=0.}
\]
The additive derivative lattice is
\[
 \boxed{
 \mathcal T_{a,b,c}
 =
 \ker(\lambda)\cap\Z^S.
 }
\]
Since the row is nonzero, its rank is $|S|-1$. The problem is not existence of additive derivatives. The problem is finding an additive derivative that remains transverse to the dependence direction introduced by $a$ and $b$.

\section{The arithmetic Wronskian}
Define
\[
 \boxed{
 W_x(a,b)
 =a\,\mathrm d_x(b)-b\,\mathrm d_x(a).
 }
\]
Writing
\[
 \mu_p=
 \begin{cases}
 -\dfrac{v_p(a)}p,&p\mid a,\\[6pt]
 \dfrac{v_p(b)}p,&p\mid b,\\[6pt]
 0,&p\mid c,
 \end{cases}
\]
we obtain
\[
 W_x(a,b)=ab\,\mu\cdot x.
\]
An additive packet is \emph{transverse} when $W_x(a,b)\neq0$. It is degenerate when the two input terms remain derivative-proportional.

Over $\R$, the additive projection is
\[
 P_{\mathrm{add}}
 =
 I-
 \frac{\lambda\lambda^{\mathsf T}}{\norm{\lambda}_2^2}.
\]
The obstruction direction inside the additive hyperplane is
\[
 \mu_\perp=P_{\mathrm{add}}\mu.
\]
When $\mu_\perp\neq0$, the rank-one obstruction projector is
\[
 P_{\mathrm{obs}}
 =
 \frac{\mu_\perp\mu_\perp^{\mathsf T}}
 {\norm{\mu_\perp}_2^2}.
\]
This is the first appearance of the projector-resolved pattern. A large ambient prime space is decomposed into one equation direction, one obstruction direction, and the remaining null packets.

\section{Radical divisibility}
Set
\[
 R=\rad(abc).
\]
The central first-order theorem is an exact divisibility law.

\begin{theorem}[Wronskian divisibility]
For every additive packet $x\in\mathcal T_{a,b,c}$,
\[
 \boxed{
 \frac{abc}{R}\mid W_x(a,b).
 }
\]
\end{theorem}

\begin{proof}
Let $p^e\parallel a$. Every term of $\mathrm d_x(a)$ is divisible by $p^{e-1}$. Since $p\nmid b$, the term $b\,\mathrm d_x(a)$ is divisible by $p^{e-1}$. Also $p^e\mid a$, so $a\,\mathrm d_x(b)$ is divisible by $p^e$, hence by $p^{e-1}$. Therefore
\[
 p^{e-1}\mid W_x(a,b).
\]
The same argument applies to primes dividing $b$.

Now let $p^e\parallel c$. Additivity gives
\[
 \mathrm d_x(c)=\mathrm d_x(a)+\mathrm d_x(b).
\]
Using $c=a+b$,
\[
 \begin{aligned}
 W_x(a,b)
 &=a\,\mathrm d_x(b)-b\,\mathrm d_x(a)\\
 &=a\bigl(\mathrm d_x(c)-\mathrm d_x(a)\bigr)-b\,\mathrm d_x(a)\\
 &=a\,\mathrm d_x(c)-c\,\mathrm d_x(a).
 \end{aligned}
\]
The first term is divisible by $p^{e-1}$ because $\mathrm d_x(c)$ is. The second is divisible by $p^e$. Hence $p^{e-1}\mid W_x(a,b)$.

The prime supports of $a$, $b$, and $c$ are disjoint. Multiplying all local divisibilities gives
\[
 \prod_{p^e\parallel abc}p^{e-1}
 =
 \frac{abc}{R}
 \mid W_x(a,b).
\]
\end{proof}

For a transverse packet,
\[
 \boxed{
 |W_x(a,b)|\ge\frac{abc}{R}.
 }
\]
This is the integer analogue of the repeated-root divisor appearing in the polynomial Mason--Stothers theorem.

\section{The master height inequality}
For every positive $n$,
\[
 \begin{aligned}
 |\mathrm d_x(n)|
 &\le
 n\norm{x}_\infty
 \sum_{p\mid n}\frac{v_p(n)}p.
 \end{aligned}
\]
Since $p\ge2$,
\[
 \frac1p\le\frac1{2\log2}\log p.
\]
Therefore
\[
 \sum_{p\mid n}\frac{v_p(n)}p
 \le
 \frac{\log n}{2\log2},
\]
and
\[
 |\mathrm d_x(n)|
 \le
 \frac{n\log n}{2\log2}\norm{x}_\infty.
\]
Consequently,
\[
 \begin{aligned}
 |W_x(a,b)|
 &\le
 a|\mathrm d_x(b)|+b|\mathrm d_x(a)|\\
 &\le
 \frac{ab}{\log2}\norm{x}_\infty\log c.
 \end{aligned}
\]
Combining this with radical divisibility gives
\[
 \boxed{
 c
 \le
 \frac{R}{\log2}\norm{x}_\infty\log c.
 }
\]

Define the shortest transverse packet
\[
 \boxed{
 \Theta(a,b,c)
 =
 \min\left\{
 \norm{x}_\infty:
 x\in\mathcal T_{a,b,c},
 W_x(a,b)\neq0
 \right\}.
 }
\]
Then
\[
 c
 \le
 \frac{R}{\log2}\Theta(a,b,c)\log c.
\]
A subpolynomial estimate
\[
 \Theta(a,b,c)=c^{o(1)}
\]
would imply the strong $abc$ inequality. Later chapters explain why this target is too close to the conjecture to be reached by generic lattice geometry alone.

\section{Three-sector compression}
Within the $a$-sector, define the integer row
\[
 \lambda^{(a)}
 =
 \left(\frac{a v_p(a)}p\right)_{p\mid a}
\]
and let
\[
 g_a=\gcd\left\{\frac{a v_p(a)}p:p\mid a\right\}.
\]
Define $g_b$ and $g_c$ similarly. A one-row Smith normal form in each sector produces a unimodular coordinate system in which only one aggregate channel survives. The additive equation becomes
\[
 \boxed{
 g_at_a+g_bt_b-g_ct_c=0.
 }
\]
The Wronskian becomes
\[
 \boxed{
 a g_bt_b-b g_at_a\neq0.
 }
\]
All remaining sector coordinates are lift gauges. They matter only because they determine the cost of realizing an aggregate channel by an integral prime packet.

For each term define
\[
 \rho_a(t)
 =
 \min\left\{
 \norm{x^{(a)}}_\infty:
 \lambda^{(a)}\cdot x^{(a)}=g_at
 \right\},
\]
with analogous functions for $b$ and $c$. Then
\[
 \Theta(a,b,c)
 =
 \min
 \max\{\rho_a(t_a),\rho_b(t_b),\rho_c(t_c)\}
\]
over aggregate solutions satisfying additivity and transversality.

This reduction separates the global two-dimensional arithmetic geometry from the local lifting cost. It is the template used repeatedly later: compress first, then retain the exact synthesis norm of the discarded coordinates.

\section{Worked example: \texorpdfstring{$3+125=128$}{3+125=128}}
The prime packets are $3$, $5$, and $2$. The additive row is
\[
 \lambda=(1,75,-448).
\]
A shortest transverse packet is
\[
 x=(-2,6,1).
\]
It gives
\[
 \mathrm d_x(3)=-2,
 \qquad
 \mathrm d_x(125)=450,
 \qquad
 \mathrm d_x(128)=448.
\]
Thus additivity holds. The Wronskian is
\[
 W_x(3,125)
 =3(450)-125(-2)
 =1600.
\]
The radical is
\[
 R=2\cdot3\cdot5=30,
\]
and
\[
 \frac{3\cdot125\cdot128}{30}=1600.
\]
Hence the radical lower bound is attained exactly.

This example becomes a recurring calibration point. Every later refinement must explain why the sharp integer divisor is visible while the corresponding height bound remains nontrivial.

\chapter{Valuation Words, Polyphase Refinement, and Local Lift Trees}

\section{Replacing one exponent by a valuation word}
The coefficient $v_p(n)$ compresses a run of valuation shells into one scalar. For a prime power $p^e$, define the valuation word
\[
 \boxed{
 V_e(T)=1+T+\cdots+T^{e-1}.
 }
\]
It obeys the concatenation identity
\[
 V_{e+f}(T)=V_e(T)+T^eV_f(T).
\]
The refined partial derivative
\[
 \partial_{p,T}(n)
 =
 \frac npV_{v_p(n)}(T)
\]
therefore satisfies the twisted product rule
\[
 \boxed{
 \partial_{p,T}(mn)
 =
 n\partial_{p,T}(m)
 +
 mT^{v_p(m)}\partial_{p,T}(n).
 }
\]
At $T=1$, the ordinary arithmetic partial derivative is recovered.

\section{Finite polyphase rings}
Fix an integer $N\ge2$ and work in
\[
 A_N=\Z[T]/(T^N-1).
\]
The residue of $V_e(T)$ records the number of valuation shells in each congruence class modulo $N$. Multiplication by $V_e$ is represented by a circulant matrix $C_{e,N}$.

\begin{theorem}[Polyphase Smith normal form]
If $\gcd(e,N)=1$, then
\[
 \boxed{
 \SNF(C_{e,N})
 =
 \operatorname{diag}(1,\ldots,1,e).
 }
\]
\end{theorem}

\begin{proof}
Let
\[
 I_N=\ker\bigl(A_N\xrightarrow{\,T\mapsto1\,}\Z\bigr)
\]
be the augmentation ideal. It is a free abelian group of rank $N-1$, and multiplication by $V_e$ preserves it. After extension to $\C$, the eigenvalues on $I_N\otimes\C$ are
\[
 V_e(\zeta)
 =
 \frac{1-\zeta^e}{1-\zeta},
 \qquad
 \zeta^N=1,
 \quad
 \zeta\neq1.
\]
Because $\gcd(e,N)=1$, the map $\zeta\mapsto\zeta^e$ permutes the nontrivial $N$th roots. Therefore
\[
 \det\bigl(V_e\mid I_N\bigr)
 =
 \prod_{\substack{\zeta^N=1\\\zeta\neq1}}
 \frac{1-\zeta^e}{1-\zeta}
 =1.
\]
Thus multiplication by $V_e$ is an automorphism of the integral lattice $I_N$. On the quotient
\[
 A_N/I_N\cong\Z,
\]
it acts as multiplication by $V_e(1)=e$. Applying the snake lemma to
\[
 0\longrightarrow I_N\longrightarrow A_N\longrightarrow\Z\longrightarrow0
\]
shows that the full map is injective and has cokernel $\Z/e\Z$. Hence its Smith normal form is
\[
 \operatorname{diag}(1,\ldots,1,e).
\]
\end{proof}

The refined block has $N-1$ integrally invertible phase channels and one aggregate multiplicity channel. Every local denominator lies in the aggregate channel.

\section{Dyadic and odd-adic primary lifts}
Factor
\[
 e=2^\alpha\prod_{\ell\text{ odd}}\ell^{\alpha_\ell}.
\]
The cokernel of the valuation block is
\[
 \Z/e\Z
 \cong
 \Z/2^\alpha\Z
 \times
 \prod_{\ell\text{ odd}}
 \Z/\ell^{\alpha_\ell}\Z.
\]
The primary factors specify the local lifting problems for the aggregate multiplicity channel. Its $2$-primary component must be treated $2$-adically, while the odd components admit the corresponding odd-adic lifts. The phase channels require no local denominator correction.

This Smith decomposition identifies where separate $2$-primary bookkeeping first enters. The additional one-level shift in multiplicative principal-unit coordinates at the prime $2$ comes from the local $2$-adic logarithm and is established later. It is not a consequence of the Smith theorem alone.

\section{Coarse and phase derivatives}
Choose a polyphase order $N$ coprime to every active valuation exponent. After the Smith decomposition, each refined prime packet has one coarse coordinate and integrally invertible phase coordinates in the common augmentation lattice. The additive equation separates accordingly.

The coarse equation recovers the original exponent-weighted row
\[
 \sum_{p\mid a}\frac{av_p(a)}p t_p
 +
 \sum_{p\mid b}\frac{bv_p(b)}p t_p
 -
 \sum_{p\mid c}\frac{cv_p(c)}p t_p
 =0.
\]
The phase equation is exponent free:
\[
 \boxed{
 \sum_{p\mid a}\frac ap y_p
 +
 \sum_{p\mid b}\frac bp y_p
 -
 \sum_{p\mid c}\frac cp y_p
 =0.
 }
\]
Define the phase derivative
\[
 D_y(n)=n\sum_{p\mid n}\frac{y_p}{p}.
\]
The same radical divisibility argument gives
\[
 \frac{abc}{R}\mid aD_y(b)-bD_y(a).
\]

Let
\[
 \sigma(a,b)
 =
 \sum_{p\mid a}\frac1p
 +
 \sum_{p\mid b}\frac1p.
\]
If $\Phi(a,b,c)$ is the shortest transverse phase packet, then
\[
 \boxed{
 c
 \le
 R\,\sigma(a,b)\,\Phi(a,b,c).
 }
\]
This can be sharper than the exponent-weighted inequality on individual triples, but finite evidence still shows positive power growth rather than a subpolynomial law.

\section{First-order no-free-refinement principle}
The shell-refined factor
\[
 \prod_{r=0}^{e-1}
 \left(1+\frac{T_r}{p}\right)
\]
has first-order expansion
\[
 1+\frac1p\sum_rT_r+O(T^2).
\]
Only the shell sum appears. Every internal phase in the kernel of shell synthesis is invisible to both first-order additivity and the first-order Wronskian.

\begin{theorem}[First-order no-free-polyphase principle]
Let
\[
 S:(T_{p,r})_{p,r}\longmapsto
 \left(\sum_rT_{p,r}\right)_p
\]
be shell synthesis. The first-order additive map and first-order Wronskian factor through $S$. Hence
\[
 \Ker S
\]
is contained in the kernels of both first-order maps.
\end{theorem}

Polyphase splitting may lower displayed child coordinates. Exact synthesis multiplies by the corresponding reconstruction norm. At first order, refinement cannot create a new obstruction direction.

\section{Second-order phase energy}
At second order,
\[
 \prod_{r=0}^{e-1}
 \left(1+\frac{T_r}{p}\right)
 =
 1+
 \frac1p\sum_rT_r+
 \frac1{p^2}\sum_{r<s}T_rT_s+
 O(T^3).
\]
Let $\widehat T_j$ be the discrete Fourier transform of the shell coordinates. Parseval gives
\[
 \sum_{r<s}T_rT_s
 =
 \frac12\left[
 \widehat T_0^2
 -
 \frac1e\sum_j
 \widehat T_j\widehat T_{-j}
 \right].
\]
For real coordinates,
\[
 \sum_{r<s}T_rT_s
 =
 \frac12\left[
 \left(1-\frac1e\right)|\widehat T_0|^2
 -
 \frac1e\sum_{j\neq0}|\widehat T_j|^2
 \right].
\]
Thus nonzero polyphase energy first becomes visible at second order. Later formal-smoothness results show that this curvature is not itself a lifting obstruction.

\section{Exact local lift gauges}
Let
\[
 r_u=\rad(u),
 \qquad
 P_u=\frac{u}{r_u}.
\]
A shell packet $y=(y_{p,r})$ produces
\[
 D_y(u)
 =
 u\sum_{p\mid u}\frac1p
 \sum_{r=0}^{v_p(u)-1}y_{p,r}
 =P_ut_u,
\]
where
\[
 t_u=
 \sum_{p\mid u}
 \frac{r_u}{p}
 \sum_r y_{p,r}.
\]
Because
\[
 \gcd_{p\mid u}\frac{r_u}{p}=1,
\]
every integer $t_u$ is attainable when $u>1$.

The shell lift gauge is
\[
 \rho_u^{\mathrm{sh}}(t)
 =
 \min\left\{
 \norm{y}_\infty:
 \sum_{p\mid u}
 \frac{r_u}{p}
 \sum_r y_{p,r}=t
 \right\}.
\]
At the aggregate level,
\[
 P_at_a+P_bt_b=P_ct_c,
\]
and
\[
 W_y(a,b)
 =
 P_aP_b(r_at_b-r_bt_a).
\]
The radical divisor is attained sharply when
\[
 r_at_b-r_bt_a=P_c.
\]
When $a,b>1$, all sharp aggregate solutions form a one-parameter family
\[
 \begin{aligned}
 t_a(n)&=t_a^{(0)}+nr_a,\\
 t_b(n)&=t_b^{(0)}+nr_b,\\
 t_c(n)&=t_c^{(0)}+nr_c.
 \end{aligned}
\]
The unit-sector cases have no moving parameter. If $a=1$, then
\[
 (t_a,t_b,t_c)=(0,P_c,P_b),
\]
and if $b=1$, then
\[
 (t_a,t_b,t_c)=(-P_c,0,-P_a).
\]
For obstruction multiple $k$, these fixed triples are multiplied by $k$. Thus the high-dimensional sharp search reduces either to one integer parameter plus three local lift gauges, or to one fixed aggregate triple in a unit-sector case.

\chapter{The Exact Obstruction Complex and Formal Shell Geometry}

\section{A joint exponent-free two-row complex}
The Smith obstruction is attached to the exponent-free phase, or shell-synthesis, rows rather than to the exponent-weighted derivative row. For every prime $p\mid abc$, define
\[
 \Lambda_p=
 \begin{cases}
 a/p,&p\mid a,\\
 b/p,&p\mid b,\\
 -c/p,&p\mid c,
 \end{cases}
 \qquad
 \mathsf w_p=
 \begin{cases}
 -ab/p,&p\mid a,\\
 ab/p,&p\mid b,\\
 0,&p\mid c.
 \end{cases}
\]
Consider
\[
 \mathsf M_{a,b,c}
 =
 \begin{pmatrix}
 \Lambda\\
 \mathsf w
 \end{pmatrix}:
 \Z^S\longrightarrow\Z^2.
\]
The first row is the exponent-free additive phase row. The second is its integer Wronskian row. Set
\[
 Q=\frac{abc}{R}.
\]

\begin{theorem}[Smith obstruction theorem]\label{thm:phase-smith}
Let $a+b=c$ be primitive, with at least two of $a,b,c$ greater than one. Then
\[
 \boxed{
 \SNF(\mathsf M_{a,b,c})
 =
 \operatorname{diag}(1,Q).
 }
\]
Consequently,
\[
 \boxed{
 \operatorname{coker}(\mathsf M_{a,b,c})_{\mathrm{tors}}
 \cong
 \Z/Q\Z.
 }
\]
The exceptional triple $(1,1,2)$ has only one nonempty prime sector, so the matrix has rank one and is excluded from the two-row statement.
\end{theorem}

\begin{proof}
The gcd of the entries of $\Lambda$ is one. Indeed, every prime dividing one sector is absent from every other nonempty sector, and an entry from another sector is not divisible by it. Hence the first Smith factor is one.

Rows restricted to one sector are proportional, so same-sector $2\times2$ minors vanish. If $p$ and $q$ lie in two different nonempty sectors, the corresponding minor is
\[
 \pm\frac{abc}{pq}
 =
 \pm Q\frac{R}{pq}.
\]
Every nonzero minor is therefore divisible by $Q$. Conversely, for each prime $\ell\mid R$, choose $p=\ell$ and choose $q$ in a different nonempty sector. Since $R$ is squarefree, $R/(pq)$ is not divisible by $\ell$. Thus the gcd of the reduced minors $R/(pq)$ is one. The gcd of all nonzero minors is exactly $Q$. The determinantal-divisor description of Smith normal form now gives the result.
\end{proof}

\section{The exact lattice extension}
Define
\[
 L=\ker(\Lambda)\cap\Z^S
\]
and
\[
 L_0=L\cap\ker(\mathsf w).
\]
Because $\Lambda$ is primitive, choose $t\in\Z^S$ with $\Lambda(t)=1$. Then
\[
 \Z^S=L\oplus\Z t.
\]
The image of $\mathsf M_{a,b,c}$ is generated by $(1,\mathsf w(t))$ together with $(0,\mathsf w(L))$. Its index in $\Z^2$ is therefore the index of $\mathsf w(L)$ in $\Z$. Theorem~\ref{thm:phase-smith} shows that this index is $Q$, so
\[
 \mathsf w(L)=Q\Z.
\]
Hence
\[
 \omega(x)=\frac{\mathsf w\cdot x}{Q}
\]
is an integer-valued surjection and gives the exact sequence
\[
 \boxed{
 0\longrightarrow L_0
 \longrightarrow L
 \xrightarrow{\omega}
 \Z
 \longrightarrow0.
 }
\]
The extension splits algebraically. The difficult quantity is the norm of a section, not its existence.

For $k\in\Z$, define
\[
 H(k)
 =
 \min\left\{
 \norm{x}_\infty:
 x\in L,
 \omega(x)=k
 \right\}.
\]
Three heights must be distinguished:
\[
 H(1),
 \qquad
 \Phi=\min_{k\neq0}H(k),
 \qquad
 \overline H=
 \inf_{k\ge1}\frac{H(k)}k.
\]
The shortest nonzero detector may live in a nonprimitive obstruction coset.

\section{The stable real height}
Relax the integral packet to a real vector in the additive hyperplane. The exact continuous width can be computed from the three sector weights. Let
\[
 \sigma_u=\sum_{p\mid u}\frac1p
\]
and
\[
 \Delta(a,b,c)
 =
 \min\left\{
 \sigma_a+\sigma_b,
 \frac cb(\sigma_a+\sigma_c),
 \frac ca(\sigma_b+\sigma_c)
 \right\}.
\]
Then
\[
 \boxed{
 \overline H
 =
 \frac{c}{R\Delta(a,b,c)}.
 }
\]
The discrete master inequality becomes
\[
 c\le R\Delta\Phi.
\]
Since
\[
 \Phi=\overline H\,k_*\Gamma
\]
for the optimal obstruction multiple $k_*$ and an integrality factor $\Gamma\ge1$,
\[
 \boxed{
 \Delta\Phi
 =
 \frac cR k_*\Gamma.
 }
\]
This is the no-escape identity. Optimizing the same first-order invariant cannot lower the continuous floor $c/R$.

\section{The full shell deformation}
Introduce one variable for every valuation shell and define
\[
 \mathscr J_n(\mathbf T)
 =
 \prod_{p^e\parallel n}
 \prod_{r=0}^{e-1}(p+T_{p,r}).
\]
At $\mathbf T=0$,
\[
 \mathscr J_n(0)=n.
\]
For $a+b=c$, set
\[
 F_{a,b,c}(\mathbf T)
 =
 \mathscr J_a(\mathbf T)
 +
 \mathscr J_b(\mathbf T)
 -
 \mathscr J_c(\mathbf T).
\]
The shell gradient has entries
\[
 \frac ap,
 \qquad
 \frac bp,
 \qquad
 -\frac cp
\]
repeated according to valuation depth. The gcd of all entries is one.

\begin{theorem}[Formal affine-space theorem]
Let $d$ be the number of shell variables. After a unimodular integral change of shell coordinates, the $(\mathbf T)$-adic completion of
\[
 \Z[\mathbf T]/(F_{a,b,c})
\]
at the zero section is isomorphic to a formal power-series ring $\Z[[U_2,\ldots,U_d]]$.
\end{theorem}

\begin{proof}
A primitive gradient row can be completed to a unimodular matrix. In the new coordinates, one variable has linear coefficient one, which is a unit in $\Z$. The formal implicit-function recursion therefore solves uniquely for that variable at every degree with integral coefficients, giving the stated isomorphism of completed rings.
\end{proof}

\begin{corollary}[Integral jet extension]
Every integral first-order tangent vector to $F_{a,b,c}=0$ extends to an integral formal curve through all orders.
\end{corollary}

The consequence is decisive. Higher shell curvature exists, but there is no formal lifting obstruction to exploit. The missing arithmetic must be a global size or conductor estimate.

\section{Gauge collapse}
Within one sector, the product
\[
 \prod_{p,r}(p+T_{p,r})
\]
depends only on the aggregate product of its shell variables. Internal phase directions that preserve this product form a gauge torus. The full shell equation factors through three aggregate sector products
\[
 A,
 \qquad
 B,
 \qquad
 C,
\]
with one linear relation
\[
 aA+bB=cC.
\]
Thus all formal complexity collapses to a two-dimensional aggregate plane plus gauge directions. This explains why local refinements repeatedly fail to create a new global obstruction.

\chapter{Exact Moving-Coset Search and the First-Order Barrier}

\section{Obstruction cosets}
Choose a section $x_0\in L$ with $\omega(x_0)=1$. Then
\[
 L
 =
 \bigsqcup_{k\in\Z}
 (kx_0+L_0).
\]
For fixed $k$, the problem is a closest-vector problem to the moving target $-kx_0$ in the lattice $L_0$. The three-sector compression makes this search much more explicit.

\section{Aggregate affine line}
Write
\[
 u=P_ur_u,
 \qquad
 r_u=\rad(u),
 \qquad
 P_u=u/r_u.
\]
The sharp aggregate conditions are
\[
 P_at_a+P_bt_b=P_ct_c
\]
and
\[
 r_at_b-r_bt_a=kP_c.
\]
Since $\gcd(r_a,r_b)=1$, choose $s,t$ with
\[
 sr_a+tr_b=1.
\]
One solution is
\[
 t_a^{(0)}=-tkP_c,
 \qquad
 t_b^{(0)}=skP_c.
\]
Every aggregate solution is
\[
 \boxed{
 \begin{aligned}
 t_a(k,n)&=t_a^{(0)}+nr_a,\\
 t_b(k,n)&=t_b^{(0)}+nr_b,\\
 t_c(k,n)&=t_c^{(0)}+nr_c,
 \end{aligned}
 \qquad n\in\Z.
 }
\]
When $a,b>1$, the global search in one obstruction coset is therefore one-dimensional. If $a=1$, the empty $a$-sector forces $t_a=0$ and fixes the coarse parameter, giving
\[
 (t_a,t_b,t_c)=(0,kP_c,kP_b).
\]
If $b=1$, the corresponding fixed triple is
\[
 (t_a,t_b,t_c)=(-kP_c,0,-kP_a).
\]
Thus a unit-sector coset has one aggregate candidate rather than a moving affine line.

\section{Local CRT normal form}
Fix one sector with primes
\[
 p_1,\ldots,p_m
\]
and radical
\[
 r=\prod_i p_i.
\]
For a nonempty sector, the local equation is
\[
 \sum_i\frac r{p_i}x_i=t.
\]
An empty unit sector has no prime coordinates and is feasible exactly when $t=0$.
Modulo $p_i$, the coordinate $x_i$ is forced:
\[
 x_i
 \equiv
 t\left(\frac r{p_i}\right)^{-1}
 \pmod{p_i}.
\]
Choose the centered residue $\xi_i$ and write
\[
 x_i=\xi_i+p_i z_i.
\]
Substitution gives one quotient equation
\[
 \sum_i z_i=q(t).
\]
For a box radius $H$, each $z_i$ lies in the interval
\[
 \ell_i(H)
 =
 \left\lceil\frac{-H-\xi_i}{p_i}\right\rceil
 \le z_i\le
 \left\lfloor\frac{H-\xi_i}{p_i}\right\rfloor
 =u_i(H).
\]
The local lift exists exactly when
\[
 \boxed{
 \sum_i\ell_i(H)
 \le q(t)\le
 \sum_i u_i(H).
 }
\]
This is an exact interval criterion.

\section{Moving-interval theorem}
For fixed $k$ and $H$, the three sectors produce three integer intervals for the coarse parameter $n$. A packet of height at most $H$ exists if and only if those intervals share an integer.

\begin{theorem}[Moving-interval theorem]
For fixed obstruction multiple $k$ and radius $H$,
\[
 H(k)\le H
\]
if and only if
\[
 I_a(H,k)
 \cap
 I_b(H,k)
 \cap
 I_c(H,k)
 \cap\Z
 \neq\varnothing.
\]
\end{theorem}

The first feasible $H$ gives the exact coset height. The preceding radius is an explicit no-smaller-box certificate.

\section{Finite global stopping}
The continuous lower bound gives
\[
 H(k)\ge k\overline H.
\]
If the best exact height found so far is $B$, then no coset with
\[
 k\overline H\ge B
\]
can improve it. Therefore only finitely many obstruction multiples must be searched:
\[
 \boxed{
 1\le k\le
 \left\lfloor\frac{B-1}{\overline H}\right\rfloor.
 }
\]
Every unsearched coset is excluded by a proved lower bound.

\section{Phase and coarse separation}
The local target has the form
\[
 t_u(k,n)=\alpha_uk+r_un.
\]
For a local prime $p\mid r_u$,
\[
 t_u(k,n)\equiv\alpha_uk\pmod p.
\]
Thus the obstruction multiple $k$ determines the finite phase packet. The coarse parameter $n$ changes only the quotient load. The exact search is a moving polyphase mechanism rather than a generic mixed-integer program.

\section{Arithmetic correction and arbitrary precision}
Some large high-quality triples exceed signed $64$-bit range, and intermediate quantities such as
\[
 ab,
 \qquad
 \frac{abc}{R},
 \qquad
 w\cdot x
\]
can overflow even when the input terms individually fit. Every public arithmetic entry point must therefore promote integer-like inputs to arbitrary-precision integers before multiplication. Floating-point values are display-only. All logical comparisons and stopping bounds are exact rational or integer computations.

Several reported ratios are mathematically at least one:
\[
 \frac{\Phi}{\overline H}\ge1,
 \qquad
 \frac{\Phi}{k_*\overline H}\ge1,
 \qquad
 \frac{H(1)}{\Phi}\ge1.
\]
Values above one are not probability errors.

\section{What the exact search proved}
The moving search found nonprimitive optima on many finite examples. It also found increasingly deep phase resonances. This showed that a fixed cutoff in $k$ is unjustified. More importantly, the no-escape identity proved that even a complete exact search of the first-order invariant cannot imply $abc$ independently. The search became a diagnostic and classification engine rather than the final proof mechanism.
